\chapter[2022 --- Paabo]{2022: Svante P\"a\"abo}
\label{ch:2022_svantepaabo}

\section*{Main Contribution}

\textit{Developed methods to sequence the genome of extinct Neanderthals and Denisovans, revealing that ancient DNA survives in modern humans and shaped human evolution.}

\section{Official Citation}

for his discoveries concerning the genomes of extinct hominins and human evolution

\section{Background and Historical Context}

Svante P\"a\"abo, a Swedish geneticist born in 1955 in Stockholm, Sweden, was awarded the 2022 Nobel Prize in Physiology or Medicine ``for his discoveries concerning the genomes of extinct hominins and human evolution.'' P\"a\"abo's pioneering work in paleogenomics---the study of ancient DNA---has fundamentally changed our understanding of human evolutionary history and has revealed that the genomes of modern humans carry the genetic legacy of interbreeding with other hominin species, including Neanderthals and Denisovans.

The study of human evolution has traditionally relied on the analysis of fossil remains---bones, teeth, and other physical artifacts left behind by our ancestors. While fossil evidence has provided invaluable information about the morphology, anatomy, and approximate timeline of human evolution, it has been limited in its ability to reveal the genetic relationships between different hominin species and the evolutionary processes that shaped the modern human genome.

By the mid-twentieth century, the understanding of human evolution was based primarily on the morphological analysis of fossil hominins. Fossils of \textit{Homo erectus}, \textit{Homo neanderthalensis}, and anatomically modern \textit{Homo sapiens} had been discovered at various sites around the world, and comparative anatomical studies had established the broad outlines of human evolutionary history. However, many fundamental questions remained unanswered. What was the genetic relationship between Neanderthals and modern humans? Had they interbred? What was the demographic history of modern human populations? Were there other hominin species that coexisted with modern humans? These questions could not be answered by morphology alone---they required access to genetic information from the ancient specimens themselves.

The idea of extracting and analyzing DNA from ancient specimens was first proposed in the 1980s, but the practical challenges were immense. DNA is a fragile molecule that degrades rapidly after death, breaking down into small fragments and accumulating chemical modifications that distort its sequence. Recovering authentic ancient DNA from specimens that are thousands or hundreds of thousands of years old seemed to many scientists to be an impossible task. Moreover, the risk of contamination with modern DNA---from the researchers handling the specimens, from the storage environment, and from other sources---posed a constant threat to the authenticity of any ancient DNA sequences that might be recovered.

P\"a\"abo was one of the first researchers to take on this challenge. Beginning in the early 1980s, when he was a graduate student at Uppsala University in Sweden, P\"a\"abo set out to develop methods for extracting and analyzing DNA from ancient specimens. His early work focused on Egyptian mummies, which provided relatively well-preserved DNA due to the dry conditions in which they were preserved. In a landmark study published in 1985, P\"a\"abo successfully extracted and characterized DNA from a 2,400-year-old Egyptian mummy, demonstrating that ancient DNA could be recovered and analyzed.

This achievement was met with skepticism by many in the scientific community, who questioned whether the DNA sequences obtained from ancient specimens were truly authentic or were the result of contamination with modern DNA. P\"a\"abo recognized that contamination was a critical concern and devoted considerable effort to developing rigorous protocols for preventing and detecting contamination. These protocols, which included clean-room procedures, negative controls, and independent replication of results, became the gold standard for ancient DNA research and remain in use today.

In the decades that followed, P\"a\"abo refined his methods and applied them to increasingly ancient and challenging specimens. He analyzed DNA from Neanderthal bones, from the permafrost-preserved remains of woolly mammoths and other Ice Age megafauna, and from the 700,000-year-old leg bone of a horse found in the Yukon territory of Canada---at the time, the oldest DNA ever sequenced. Each new technical achievement pushed the boundaries of what was possible in ancient DNA research and expanded the temporal and geographic range of specimens that could be studied.

P\"a\"abo's development of next-generation sequencing technologies for ancient DNA was particularly transformative. Traditional Sanger sequencing methods, which had been used for the initial ancient DNA studies, were slow, labor-intensive, and required relatively large amounts of starting material. P\"a\"abo recognized that the emerging massively parallel sequencing technologies---which could sequence millions of DNA fragments simultaneously---were ideally suited for ancient DNA work, where the starting material is highly fragmented and present in tiny quantities. He adapted these technologies for ancient DNA analysis, developing new computational methods for distinguishing authentic ancient DNA sequences from contamination and sequencing errors.

\section{The Discovery/Contribution}

P\"a\"abo's most significant contribution to the understanding of human evolution was the sequencing of the Neanderthal genome. Neanderthals (\textit{Homo neanderthalensis}) were a closely related hominin species that lived in Europe and western Asia from approximately 400,000 to 40,000 years ago. They were adapted to cold climates, had robust bodies and large brains, and coexisted with modern humans (\textit{Homo sapiens}) for thousands of years before going extinct. The question of whether modern humans and Neanderthals interbred had been debated for decades, but could not be resolved using fossil evidence alone.

P\"a\"abo's approach to sequencing the Neanderthal genome was based on the development of new technologies for working with highly degraded ancient DNA. He and his colleagues developed methods for extracting and purifying DNA fragments from Neanderthal bone specimens, constructing DNA libraries from these fragments, and sequencing them using high-throughput sequencing technologies. The resulting sequences were then aligned to the human genome reference sequence to identify Neanderthal-specific variants.

The sequencing of the Neanderthal genome, which P\"a\"abo and his colleagues published in a landmark paper in 2010, was a monumental technical achievement. The genome was assembled from DNA extracted from bone specimens of three Neanderthal individuals found in Vindija Cave, Croatia, dating to approximately 40,000 years ago. The quality of the Neanderthal genome sequence was remarkably high, with coverage comparable to that of early drafts of the human genome.

The analysis of the Neanderthal genome revealed a striking finding: modern non-African humans carry approximately 1--4\% of their genome from Neanderthal ancestors. This finding provided definitive evidence that modern humans and Neanderthals had interbred during their coexistence in Europe and western Asia, approximately 50,000--60,000 years ago. The Neanderthal DNA segments that persist in modern human genomes are found at frequencies consistent with having been inherited from a small number of interbreeding events and subsequently maintained by natural selection.

P\"a\"abo's work also led to the discovery of a previously unknown hominin species---the Denisovans. In 2010, P\"a\"abo and his colleagues analyzed DNA extracted from a small finger bone found in Denisova Cave in Siberia, Russia. The DNA sequence was unlike that of any known hominin species, including both modern humans and Neanderthals, leading to the identification of a distinct population of archaic humans that was named the Denisovans, after the cave where the specimen was found.

The analysis of Denisovan DNA revealed that Denisovans were closely related to Neanderthals but genetically distinct from them. Like Neanderthals, Denisovans interbred with modern humans, and Denisovan DNA is found in the genomes of modern Melanesian and Southeast Asian populations at frequencies of up to 5--6\%. The discovery of Denisovans demonstrated that the human evolutionary tree was far more complex and bushy than previously appreciated, with multiple hominin species coexisting and interacting in different parts of the world.

P\"a\"abo's work also revealed the functional significance of archaic human DNA in modern human genomes. Neanderthal and Denisovan DNA segments have been shown to influence a wide range of traits in modern humans, including immune function, skin pigmentation, hair morphology, and susceptibility to disease. For example, a Neanderthal gene variant that affects the immune system's response to viral infections is carried by approximately 50\% of modern Eurasians, suggesting that it provided a selective advantage to modern humans who acquired it through interbreeding with Neanderthals. The gene in question, \textit{TLR1}, encodes a Toll-like receptor that recognizes bacterial lipopeptides, and the Neanderthal variant enhances the innate immune response to bacterial infections.

The study of archaic human DNA has also revealed that different human populations carry different amounts of Neanderthal and Denisovan DNA, reflecting the complex history of interbreeding between modern humans and archaic hominins. European and Asian populations carry approximately 1--4\% Neanderthal DNA, while Melanesian and Southeast Asian populations carry up to 5--6\% Denisovan DNA. Some African populations, once thought to carry no archaic DNA, have been shown to carry small amounts of Neanderthal and Denisovan genetic material, suggesting more complex patterns of admixture than initially appreciated.

Conversely, some Neanderthal DNA segments appear to have been deleterious to modern humans and have been progressively eliminated from the genome by natural selection. Regions of the modern human genome that are depleted of Neanderthal DNA include genes involved in brain development and male fertility, suggesting that Neanderthal variants in these regions were harmful and were selected against. The X chromosome and regions near genes expressed in the testes are particularly depleted of archaic DNA, consistent with the hypothesis that hybrid male fertility was reduced---a common feature of interspecific hybridization in many animal species.

\section{Impact on Modern Medicine}

P\"a\"abo's discoveries have had a significant impact on medicine, particularly in the areas of disease genetics, immune function, and personalized medicine.

\textbf{Disease susceptibility:} The analysis of archaic human DNA has revealed that Neanderthal and Denisovan genetic variants influence susceptibility to a range of modern diseases. For example, a Neanderthal gene variant has been associated with increased risk of severe COVID-19, while another variant has been associated with reduced risk. These findings have provided insights into the genetic basis of disease susceptibility and have highlighted the role of evolutionary history in shaping modern human health.

\textbf{Immune function:} Many Neanderthal gene variants that have been retained in modern human genomes are involved in immune function, suggesting that interbreeding with Neanderthals provided modern humans with genetic adaptations that enhanced their immune defenses. The understanding of how archaic human DNA influences immune function has informed the development of strategies for combating infectious diseases and has provided insights into the mechanisms of immune regulation.

\textbf{Skin and hair biology:} Neanderthal DNA variants have been associated with differences in skin pigmentation, hair morphology, and other integumentary traits. These variants may have provided adaptive advantages to modern humans who migrated to different environments, such as the reduced UV radiation at higher latitudes. The understanding of how archaic human DNA influences skin and hair biology has implications for dermatology and the study of human adaptation to different environments.

\textbf{Human genetics and personalized medicine:} The study of archaic human DNA has highlighted the importance of evolutionary history in shaping modern human genetic diversity and has provided a framework for understanding how genetic variants influence health and disease. This understanding has implications for personalized medicine, as the effects of genetic variants on disease risk and drug response may be influenced by their evolutionary origin.

Svante P\"a\"abo was awarded the Nobel Prize in Physiology or Medicine in 2022 ``for his discoveries concerning the genomes of extinct hominins and human evolution.'' His work has fundamentally changed our understanding of human evolutionary history and has provided new insights into the relationship between evolution and modern human health.

\section{Legacy}

Svante P\"a\"abo's legacy extends across multiple areas of biology and medicine. His pioneering work in ancient DNA research has created a new field---paleogenomics---that has transformed our understanding of human evolution and has provided new tools for studying the genetic history of populations and species.

The development of methods for extracting and analyzing ancient DNA has enabled the study of evolutionary processes on timescales that were previously inaccessible to genetic analysis. Ancient DNA studies have revealed the complex population dynamics of human evolution, including multiple waves of migration, admixture between different hominin populations, and the role of natural selection in shaping the modern human genome.

P\"a\"abo's work has also established the Max Planck Institute for Evolutionary Anthropology in Leipzig, Germany, as one of the world's leading centers for ancient DNA research. The institute's paleogenomics laboratory has produced genome sequences from dozens of ancient hominin specimens, providing an increasingly detailed picture of human evolutionary history.

The field of ancient DNA research continues to expand rapidly, with new technologies enabling the recovery of DNA from increasingly ancient specimens and from a wider range of organisms. The legacy of P\"a\"abo's work provides the foundation for these ongoing investigations and will continue to shape our understanding of human evolution and its impact on modern health for generations to come.


\section{Modern Developments}

Pääbo's paleogenomics work has led to modern human evolutionary genetics. Ancient DNA has revealed interbreeding between modern humans and Neanderthals, with ~2% of non-African genomes derived from Neanderthals. Understanding of ancient DNA has revealed adaptive introgression---genes inherited from Neanderthals that help modern humans (e.g., immune genes, high-altitude adaptation). Understanding of Neanderthal genomes has informed research on COVID-19 susceptibility (Neanderthal haplotype affects severe disease risk). Paleogenomics continues to reveal human evolutionary history.


\section{Bibliography}

\begin{thebibliography}{9}
\bibitem{nobel-2022}
Nobel Prize Outreach, ``The Nobel Prize in Physiology or Medicine 2022,''\emph{NobelPrize.org}.\\
\url{https://www.nobelprize.org/prizes/medicine/2022/summary/}

\bibitem{lecture-2022}
Svante P"a"abo, ``Nobel Lecture,''\emph{NobelPrize.org}. Winner-authored primary source; downloadable from the official Nobel Prize archive.\\
\url{https://www.nobelprize.org/prizes/medicine/2022/paabo/lecture/}
\end{thebibliography}
